\documentclass[12pt,twoside=off,
bibtotoc,liststotoc,appendixprefix,paper=a4,headings=small]{scrbook} % 'twoside=on' für Druckversion oder 'twoside=off' für Onlineversion

\makeatletter
\renewcommand\chapter{\thispagestyle{plain}%
\global\@topnum\z@
\@afterindentfalse
\secdef\@chapter\@schapter}
\makeatother

%
% Packages
% -----------------------------------
\usepackage[
  paper=a4paper,
  left=12.5mm,
  right=25mm,
  top=25mm,
  bottom=50mm,
  bindingoffset=10mm]{geometry}		% Seitenränder und Bindungskorrektur einstellen
  
\usepackage{apacite} 				% Literatur-Referenzen: American Psycholog. Assoc.
\usepackage{natbib}					

\setcitestyle{round,aysep={}} 		% Indizierg. in runden Klammern, zw. Autor u. Jahr
\usepackage[utf8]{inputenc}         % Umlaute im Text
\usepackage{ngerman}				% Rechtschreibg.
\usepackage[T1]{fontenc}
\usepackage{lmodern}				% Schriftfamilie
\usepackage{microtype}				% für die Mikrotypografie (besseres Schriftbild)
\setkomafont{disposition}{\bfseries}

\usepackage{graphicx} 				% Grafiken einfügen (pdf,png - aber jpg vermeiden)
\graphicspath{{./Bilder/}}          % Pfad zu den Bildern

\usepackage{url}					% URL's formatieren (z.B. in Literatur) 
\usepackage[colorlinks,linkcolor=black,citecolor=black,urlcolor=black]{hyperref} 				% für Hyperlinks in PDF-Dokumenten   
  
\usepackage{tabularx} 				% bessere Gestaltung von Tabellen
\usepackage{longtable} 		
\usepackage{multicol}				
\usepackage{multirow}
\usepackage{booktabs}
\usepackage{tabularx}
		
\usepackage[active]{srcltx}

\usepackage{listings}				% Algorithmen

\usepackage{mdwlist}				% Listen

\usepackage{setspace} 				% Zeileneinstellung
\newtheorem{mydef}{Merksatz}  		% Falls Beispiele, Merksätze m. fortl. Nr. gebr. werden
\newtheorem{bsp}{Beispiel}

\usepackage{todonotes}				% zum Erstellen von ToDos im Editor

\usepackage{lscape}					% zum Rotieren von Seiten

\usepackage{amsmath}				% zum Schreiben von mathematischen Formeln

\usepackage{calc}

\usepackage{footnote}				% Fußnoten
\usepackage{tablefootnote}			% Fußnoten in Tabellen

%\clubpenalty = 10000
%\widowpenalty = 10000 \displaywidowpenalty = 10000

\hyphenation{voll-st\"andigen}		% Worttrennungen global definieren

\setcounter{tocdepth}{1}			% Ebenen, die im Inhaltsverzeichnis angezeigt werden

% Document
% -----------------------------------
\begin{document}

\frontmatter 
    % Titelseite soll keine Kopf oder Fußzeile haben
\thispagestyle{empty}

% Alle Elemente sollen zentriert sein
\begin{center}

\vspace*{-10mm}

{\LARGE FACHBEREICH \\WIRTSCHAFTSWISSENSCHAFTEN\\[1mm]}
DER FREIEN UNIVERSITÄT BERLIN\\

\vspace*{1cm}

\includegraphics[width=0.18\textwidth]{fu_logo}

\vspace*{1cm}

% Art der Arbeit => (Bachelorarbeit ,Diplomarbeit, Masterarbeit, Seminararbeit)
{\Large \textbf{Seminararbeit}}\\ 

\vspace{1cm}

% Titel der Arbeit 
{\Large \textbf{The Economic Relationship}}\\ 
\vspace*{1mm}
{\Large \textbf{between China and Russia}}\\ 
\vspace*{1mm}
{\Large \textbf{before and after the Russia-Ukraine War}}\\

\vspace{1.5cm}

% Name des/der Autors/Autoren
{\LARGE Katharina Hansen}\\[15mm]

% Gutachter, Kontaktdaten und Abgabetermin
\parbox{120mm}{
\begin{large}
\begin{tabbing}
Gutachter: \hspace{.7cm} \=Prof. Dr. Laike Yang\\[4mm]
Semester:\> Sommersemester 2025\\
Verfasser:\> Katharina Hansen\\ % alphabetische Reihenfolge (Nachname)
Matrikel-Nr.:\> 5576002\\
Email:\> katharih04@zedat.fu-berlin.de\\
Studienfach:\>Bachelor VWL\\[8mm]
\end{tabbing}
\end{large}
}

\end{center}
\clearpage{\pagestyle{empty}\cleardoublepage}
 			% Titelblatt
    \clearpage{\pagestyle{empty}\cleardoublepage}
    \onehalfspacing                  	% Zeilenabstand ab hier 1.5 zeilig
    \tableofcontents 					% Inhaltsverzeichnis
    \clearpage{\pagestyle{empty}\cleardoublepage} 
% -----------------------------------
\mainmatter 							% die einzelnen Kapitel
\chapter{Einleitung}
\chapter{Hintergrund}

% -----------------------------------
\backmatter 
\bibliographystyle{apacite}				% bei natbib in deutsch
\bibliography{./Literatur/quellen}		% Literaturquellen einbinden 

\end{document}